% 基本的な数学パッケージ
\usepackage{amsmath,amssymb}

% ブラケット記法 (\bra, \ket, \braket 等) を直感的に書くためのパッケージ
\usepackage{physics}

% 太字ベクトル用
\usepackage{bm}

% 虚数単位、自然対数の底、微分記号を立体にする
\newcommand{\ii}{\mathrm{i}}
\newcommand{\ee}{\mathrm{e}}
\RenewDocumentCommand{\dd}{}{\mathrm{d}}

% 量子回路図描画用パッケージ
\usepackage{graphicx} % \resizebox などで使用
\usepackage{tikz}
\usetikzlibrary{quantikz2} % 最新環境ではquantikz2が推奨されます
% 格子配置図・時空図・ブロック図で使う TikZ ライブラリ
\usetikzlibrary{positioning,calc,arrows.meta,fit,backgrounds,patterns,decorations.pathreplacing,shapes.geometric,3d}
% 図版の共通スタイル
\tikzset{
  xtype/.style={red!70!black},          % X 型（primal）
  ztype/.style={blue!70!black},         % Z 型（dual）
  ytype/.style={green!50!black},        % Y 型
  dataqubit/.style={circle, draw, fill=white, minimum size=4pt, inner sep=0pt},
  ancqubit/.style={circle, draw, fill=black!45, minimum size=4pt, inner sep=0pt},
  logpatch/.style={draw, dashed, thick}, % 論理アンシラパッチ
  xplaq/.style={fill=red!12, draw=red!70!black},
  zplaq/.style={fill=blue!12, draw=blue!70!black},
}

% ハイパーリンクの設定
\usepackage{xurl}
\usepackage[unicode,hidelinks]{hyperref}

% 長いコードパスを行末で安全に折り返す
\newcommand{\codepath}[1]{\nolinkurl{#1}}
\newcommand{\codeinline}[1]{\nolinkurl{#1}}

% 長い日本語文とインラインコードが混在する段落の微小なはみ出しを防ぐ
\setlength{\emergencystretch}{2em}

% 本文幅に追従する比較表
\usepackage{tabularx}

% ソースコード表示用の設定 (Python/Qiskit等を想定)
\usepackage{listings}
\usepackage{xcolor}
% listingsのLOL出力を無効化（\@writefileのバッファ問題回避）
\makeatletter
\AtBeginDocument{\renewcommand{\lstlistoflistings}{}}
\AtEndDocument{%
  \let\orig@end@writefile\@writefile
  \renewcommand{\@writefile}[2]{%
    \def\@end@tmp@ext{lol}%
    \def\@end@tmp@arg{#1}%
    \ifx\@end@tmp@ext\@end@tmp@arg
    \else
      \orig@end@writefile{#1}{#2}%
    \fi
  }%
}
\makeatother

% コールアウト用ボックス
\usepackage{tcolorbox}
\tcbuselibrary{breakable}

% チャプターゴール / Note / 警告 / 演習 環境
\newtcolorbox{chaptergoal}{
  colback=blue!5,
  colframe=blue!60!black,
  title=\textbf{この章で学ぶこと},
  fonttitle=\bfseries,
  breakable
}
\newcommand{\cgitem}{\textbullet\ }

\newtcolorbox{notebox}{
  colback=blue!5,
  colframe=blue!60!black,
  title=\textbf{Note},
  fonttitle=\bfseries,
  breakable
}

\newtcolorbox{warningbox}{
  colback=yellow!5,
  colframe=orange,
  title=\textbf{注意},
  fonttitle=\bfseries,
  breakable
}

\newtcolorbox{exercise}[1][演習]{
  colback=green!5,
  colframe=green!50!black,
  title=\textbf{#1},
  fonttitle=\bfseries,
  breakable
}

\lstset{
    language=Python,
    basicstyle=\ttfamily\small,
    keywordstyle=\color{blue}\bfseries,
    stringstyle=\color{red},
    commentstyle=\color{green!50!black},
    numbers=left,
    numberstyle=\tiny\color{gray},
    stepnumber=1,
    numbersep=5pt,
    backgroundcolor=\color{black!5},
    frame=single,
    rulecolor=\color{black!20},
    breaklines=true,
    tabsize=4,
    captionpos=b,
    showstringspaces=false
}
